\section{Why Melanoma Detection Is a Timing Problem}

Melanoma is survivable when it is caught early and frequently fatal when it is not.
Five-year survival sits near 100\% for lesions found while still localised and drops
below 30\% once the disease has metastasised~\cite{esteva2017dermatologist}. Nothing
about the tumour changes in that interval except how long it has been left alone. The
clinical problem is therefore less about treatment than about how quickly a suspicious
lesion reaches somebody qualified to look at it.

That is a resource problem. Dermatologists are unevenly distributed, waiting lists are
long, and the people most at risk are often the least likely to present early. A triage
tool that runs on a phone, in a general practice, or in a rural clinic without
specialist cover does not need to replace a dermatologist. It needs to sort lesions well
enough to move the worrying ones up the queue.

Deep learning has been able to do this in principle since Esteva et
al.~\cite{esteva2017dermatologist} trained an Inception-v3 network on roughly 129{,}000
clinical and dermoscopic images and matched board-certified dermatologists on biopsy-proven
cases. The result has been reproduced in various forms often enough that the open
question is no longer whether a convolutional network can classify skin lesions. It is
which network you should actually deploy, and what that choice costs you.

\section{The Gap Between Benchmark and Phone}

Published accuracy is usually reported for architectures that were never meant to run
on a handset. ResNet-50~\cite{he2016resnet} is the standard medical-imaging baseline and
serialises to 94.3~MB in the seven-class configuration used in this thesis.
MobileNetV2~\cite{sandler2018mobilenetv2}, designed for exactly the on-device case,
comes to 9.1~MB. Both figures are measured here rather than quoted, and the ratio is
close to eleven to one.

An eleven-fold difference in size is not a rounding error. It decides whether a model
ships inside an app or sits behind a network call, and a network call is precisely what
an offline screening tool cannot depend on. So the practical question is narrow and
answerable: what does the smaller network actually give up?

The published record answers this badly, for a reason worth stating plainly. Studies
comparing architectures on HAM10000 tend to vary more than the architecture. Split
protocol, augmentation, optimiser, batch size and epoch budget move together from paper
to paper, so a reported gap between two backbones conflates the backbone with everything
else that changed. Headline accuracy makes it worse. HAM10000 is 66.9\% benign
melanocytic nevi, so a model that answers \emph{nevus} to every image scores 66.9\%
without ever recognising a melanoma. That number looks respectable next to much of the
literature, and it is worthless.

\section{What This Thesis Does}

This thesis runs the comparison as a controlled experiment. One pipeline, one split, one
optimiser, one schedule, and exactly two things allowed to vary: the architecture and
the strategy used to handle class imbalance. That gives a $2\times3$ factorial over
MobileNetV2 and ResNet-50 crossed with augmentation-only, weighted cross-entropy and
random oversampling.

A third architecture, EfficientNet-B3~\cite{manole2024efficientnet}, was added after the
original six runs were designed and executed. It is treated throughout as an ablation
rather than as a third arm of the confirmatory design, and it is reported in its own
statistical family. The distinction matters: the first six runs were specified before
any results were seen, and folding a later addition into the same analysis would let a
post-hoc decision change whether the pre-registered findings hold. EfficientNet-B3 is run
twice, once at 224 pixels to match the other two backbones and once at its native 300,
so that any difference can be attributed to the architecture rather than to the larger
input it normally receives.

Performance is reported per class, with melanoma and basal cell carcinoma as the
headline, because those are the classes a screening tool is answerable for. Overall
accuracy appears in the results tables for completeness and is never used to decide
anything.

\section{Contributions}

\begin{enumerate}
  \item A controlled factorial on HAM10000 in which architecture and imbalance strategy
        are isolated from one another and from the rest of the pipeline, with the split
        hash recorded in every run so that any deviation is detectable rather than
        assumed absent.
  \item Per-class precision, recall and F1 for all seven classes across every
        configuration, rather than the single accuracy figure that dominates the
        published record.
  \item An EfficientNet-B3 ablation at two input resolutions, which separates the
        question of whether the architecture is better from the question of whether it
        was simply given more pixels.
  \item Paired McNemar tests with Holm-Bonferroni correction, kept in two families so
        that pre-registered comparisons and later additions are not mixed.
  \item Measured efficiency figures (parameter count, on-disk size, CPU latency) rather
        than the estimates usually quoted from the original architecture papers.
  \item A reproducible implementation in which every reported number is generated from
        the stored run artefacts, so that the tables in this document cannot drift from
        the experiments that produced them.
\end{enumerate}

\section{Structure}

Chapter~2 reviews the relevant literature and identifies what the existing record does
and does not settle. Chapter~3 states the research gap and the three research questions,
each with a pass condition fixed in advance. Chapter~4 describes HAM10000 and the
lesion-level split protocol. Chapter~5 gives the methodology, including the code that
was actually executed. Chapter~6 sets out the evaluation method and the statistical
procedure. Chapter~7 reports results, Chapter~8 discusses them, and Chapter~9 concludes.
